\section*{Bab \ref{ch:g6:propdata} --- Perbandingan Senilai dan Data}

\begin{solution}{exo:g6:propdata:1}
Apel: \omterm{def:g3:division:remainder}{hasil baginya} $\frac52 = 2.5$, $\frac{7.5}{3} = 2.5$,
$\frac{12.5}{5} = 2.5$ --- semuanya sama: \omterm{def:g6:propdata:def}{sebanding} (harganya
$2.50$ per kg).

Umur dan tinggi: $\frac{85}{2} = 42.5$ tetapi $\frac{103}{4} = 25.75$
--- tidak sama: tidak \omterm{def:g6:propdata:def}{sebanding}.
\end{solution}

\begin{solution}{exo:g6:propdata:2}
Satu croissant: $4.80 \div 4 = 1.20$ euro. Tujuh:
$7 \times 1.20 = 8.40$ euro.
\end{solution}

\begin{solution}{exo:g6:propdata:3}
Koefisiennya: $18 \div 10 = 1.8$ euro per liter. Lalu
$25$ L berharga $25 \times 1.8 = 45$ euro; $72$ euro membeli
$72 \div 1.8 = 40$ L; $60$ L berharga $60 \times 1.8 = 108$ euro.
\end{solution}

\begin{solution}{exo:g6:propdata:4}
Untuk $250$ km: $2.5$ kali bahan bakar untuk $100$ km, jadi
$2.5 \times 6 = 15$ L. Untuk $350$ km: $3.5 \times 6 = 21$ L. Dengan
$27$ L: $27 \div 6 = 4.5$ ratusan km, yaitu $450$ km.
\end{solution}

\begin{solution}{exo:g6:propdata:5}
$50\,\%$ dari $84$: setengahnya, $42$. \quad
$25\,\%$ dari $200$: \omterm{def:g3:division:half}{seperempatnya}, $50$. \quad
$10\,\%$ dari $63$: sepersepuluhnya, $6.3$. \quad
$30\,\%$ dari $70$: $3 \times 7 = 21$.
\end{solution}

\begin{solution}{exo:g6:propdata:6}
$60\,\%$ dari $450$: $\frac{60}{100} \times 450 = 270$ murid makan di
kantin; $450 - 270 = 180$ murid tidak.
\end{solution}

\begin{solution}{exo:g6:propdata:7}
Kurang dari $10$ buku: Selasa ($8$) dan Kamis ($6$). Rabu dikurangi
Kamis: $17 - 6 = 11$ buku lebih banyak.
\end{solution}

\begin{solution}{exo:g6:propdata:8}
Diagram batang dengan lima batang setinggi $14$, $16$, $13$, $17$,
$20$ (skala setiap $2$ atau $5$ derajat berjalan baik). Hari paling
panas: Jumat ($20$ derajat).
\end{solution}

\begin{solution}{exo:g6:propdata:9}
Untuk $10$ orang, kalikan dengan $\frac{10}{6} = \frac53$: tepungnya
$450 \times \frac{10}{6} = 750$ g. Telurnya:
$3 \times \frac{10}{6} = 5$ --- untunglah bilangan bulat. (Kalau
bukan, orang akan \omterm{def:g4:numbers:round}{membulatkannya} \emph{ke atas} agar cukup.)
\end{solution}

\begin{solution}{exo:g6:propdata:10}
\emph{1.} Dalam $2$ jam: $48$ km. Dalam setengah jam: $12$ km. Dalam
$1$ jam $30$: $24 + 12 = 36$ km.

\emph{2.} $60 \div 24 = 2.5$ jam, yaitu $2$ jam $30$ menit.
\end{solution}

\begin{solution}{exo:g6:propdata:11}
Bola: potongannya $20\,\%$ dari $15 = 3$ euro, harga barunya $12$
euro. Raket: potongannya $8$ euro, harga barunya $32$ euro. Harga
barunya adalah $80\,\%$ harga lamanya pada setiap keadaan: \omterm{def:g6:propdata:def}{sebanding},
dengan koefisien $0.8$.
\end{solution}

\begin{solution}{exo:g6:propdata:12}
Setelah $t$ jam, lilin yang tebal sudah terbakar $\frac t6$ tingginya:
ia berdiri setinggi $1 - \frac t6$ dari tinggi awalnya; yang tipis
setinggi $1 - \frac t4$. Kita menginginkan
\[
1 - \frac t4 = \frac12 \left(1 - \frac t6\right)
\quad\text{yaitu}\quad
1 - \frac t4 = \frac12 - \frac t{12}.
\]
Lalu $\frac12 = \frac t4 - \frac t{12} = \frac{3t - t}{12} =
\frac{t}{6}$, jadi $t = 3$ jam. Periksa: setelah $3$ jam lilin yang
tebal berdiri setinggi $1 - \frac36 = \frac12$ dan yang tipis setinggi
$1 - \frac34 = \frac14$ --- memang setengahnya.
\end{solution}

\begin{solution}{pb:g6:propdata:1}
\textbf{1.} $100\,000$ cm $= 1\,000$ m $= 1$ km: satu sentimeter
pada peta adalah satu kilometer di tanah.

\textbf{2.} $7.5$ cm $\rightarrow 7.5$ km antara kedua desa;
$12$ cm $\rightarrow 12$ km tepi danau.

\textbf{3.} $20$ km $\rightarrow 20$ cm pada peta.

\textbf{4.} $5$ km $= 500\,000$ cm, jadi skalanya
$1:500\,000$. Peta pendakiannya ($1:100\,000$) yang lebih rinci:
ia memakai $5$ cm kertas ketika peta yang lain hanya
menghabiskan $1$ cm.

\textbf{5.}
\begin{center}
\renewcommand{\arraystretch}{1.2}
\begin{tabular}{l|ccc}
peta (cm) & $1$ & $2$ & $7.5$ \\
\hline
sebenarnya (km) & $1$ & $2$ & $7.5$ \\
\end{tabular}
\end{center}
Jarak sebenarnya $=$ jarak pada peta $\times 1$ (dalam satuan ini):
selalu pengali yang \emph{sama}, dan itu tepat definisi
\omterm{def:g6:propdata:def}{perbandingan senilai}
(\cref{def:g6:propdata:def}). Koefisiennya di sini $1$
kilometer per sentimeter.

\textbf{6.} $4$ m $= 400$ cm dan $400 \div 50 = 8$;
$3$ m $= 300$ cm dan $300 \div 50 = 6$: denahnya menunjukkan
persegi panjang $8$ cm $\times$ $6$ cm.

\textbf{7.} Tempat tidur: $200 \div 50 = 4$ cm kali $90 \div 50 = 1.8$ cm.
Pintu: $80 \div 50 = 1.6$ cm.

\textbf{8.} Luas sebenarnya: $4 \times 3 = 12$ m$^2$. Luas denahnya:
$8 \times 6 = 48$ cm$^2$. Setelah diubah:
$12$ m$^2 = 12 \times 10\,000 = 120\,000$ cm$^2$, dan
\[
120\,000 \div 48 = 2\,500 .
\]
\omterm{def:g3:division:remainder}{Hasil baginya} bukan $50$ melainkan $2\,500 = 50 \times 50$.

\textbf{9.} Satu cm$^2$ kertas mewakili persegi sebenarnya berukuran
$50$ cm kali $50$ cm, yang memuat
$50 \times 50 = 2\,500$ cm$^2$. Jadi ketika panjang dibagi
$50$, luas dibagi $50 \times 50 = 2\,500$: luas mengikuti
\emph{kuadrat} skalanya.

\textbf{10.} $6 \times 2\,500 = 15\,000$ cm$^2$, dan
$15\,000 \div 10\,000 = 1.5$ m$^2$.

\textbf{11.} Gambarnya menyarankan bahwa hari Minggu menjual
\emph{dua kali} lebih banyak (batangnya dua kali lebih tinggi). Yang
sebenarnya: kenaikannya $5$ euro dari $105$, dan
$5 \div 105 \approx 0.048$: kira-kira $5\,\%$ lebih banyak, bukan
$100\,\%$ lebih banyak. Manajernya mengabaikan peringatan untuk
memeriksa bahwa sumbunya bermula dari $0$
(\cref{met:g6:propdata:read}).

\textbf{12.} Dengan sumbu yang bermula dari $0$, batangnya naik sampai
$105$ dan $110$ satuan kecil: dua batang yang tingginya hampir sama,
yaitu gambar yang jujur untuk \omterm{ex:g1:subtraction:difference}{selisih} $5\,\%$.

\textbf{13.} Naik: kenaikannya $20$ perangko dari awal $40$:
$\frac{20}{40} = 50\,\%$. Turun: penurunannya $20$ perangko dari awal
$60$: $\frac{20}{60} = \frac13 \approx 33\,\%$. Sama-sama $20$
perangko, tetapi nilai awalnya berbeda --- persen selalu berupa
\omterm{def:g6:fractions:def}{pecahan} \emph{dari sesuatu}, dan ``sesuatu'' itu sudah berubah.

\textbf{14.} Setelah potongan $20\,\%$: $100 - 20 = 80$ euro.
Tambahan $10\,\%$ berlaku pada $80$: potongannya $8$ euro, harga
akhirnya $72$ euro. Potongan seluruhnya: $28$ euro dari $100$, jadi
$28\,\%$ --- bukan $30\,\%$. Potongan kedua bekerja pada harga yang
sudah dipotong, jadi euronya lebih kecil: persen dari besaran yang
berbeda tidak dapat dijumlahkan.

\textbf{15.} Pada peta itu, $1$ cm mewakili $5$ km, jadi $1$
cm$^2$ mewakili $5 \times 5 = 25$ km$^2$ (aturan pertanyaan 9).
Hutannya: $3 \times 25 = 75$ km$^2$. Panjang mengikuti
skalanya; luas mengikuti \emph{kuadratnya}.

\end{solution}
